
White-box attacks assume full access to the model's 
architecture, parameters, and gradients. Unlike black-box 
scenarios where the attacker has no knowledge of the 
internal structure, white-box attacks exploit the 
transparency of the model to systematically craft 
adversarial examples. Formally, given a model $f_\theta$, 
an input image $x$ with true label $y$, and a loss 
function $\mathcal{L}$, the adversarial attack objective 
is to find a perturbation $\delta$ such that:

\begin{equation}
x_{adv} = x + \delta, \quad 
\text{subject to} \quad \|\delta\|_p \leq \epsilon
\label{eq:advobj}
\end{equation}

Where $\epsilon$ controls the magnitude of the 
perturbation and $\|\cdot\|_p$ denotes the $\ell_p$ norm. 
The attacker aims to maximize the loss while keeping 
the perturbation imperceptible:

\begin{equation}
\delta^* = \arg\max_{\|\delta\|_p \leq \epsilon} 
\mathcal{L}(f_\theta(x + \delta), y)
\label{eq:advmax}
\end{equation}

\subsubsection{Fast Gradient Sign Method (FGSM)}

The Fast Gradient Sign Method (FGSM), introduced by 
Goodfellow et al. (2015), generates adversarial examples 
by computing a single-step perturbation along the sign 
of the gradient of the loss with respect to the input:

\begin{equation}
x_{adv} = x + \epsilon \cdot 
\text{sign}\left(\nabla_x 
\mathcal{L}(f_\theta(x), y)\right)
\label{eq:fgsm}
\end{equation}

Where $\nabla_x \mathcal{L}(f_\theta(x), y)$ is the 
gradient of the loss with respect to the input image 
$x$, and $\text{sign}(\cdot)$ extracts the sign of 
each gradient component. FGSM exploits the linearity 
of deep neural networks — components such as ReLU, 
defined in Equation \ref{eq:relufunc}, are intentionally 
designed to behave linearly, causing small perturbations 
applied consistently across input dimensions to produce 
drastic changes in the model output (Goodfellow et al., 
2015). The connection between memorization and FGSM 
vulnerability is examined in Section 4.

\subsubsection{Projected Gradient Descent (PGD)}

Projected Gradient Descent (PGD), proposed by Madry 
et al. (2018), extends FGSM by applying multiple 
gradient steps with projection back onto the 
$\epsilon$-ball at each iteration:

\begin{equation}
x_{adv}^{(t+1)} = \Pi_{x+\mathcal{S}} \left( 
x_{adv}^{(t)} + \alpha \cdot \text{sign}\left(
\nabla_x \mathcal{L}(f_\theta(x_{adv}^{(t)}), y)
\right)\right)
\label{eq:pgd}
\end{equation}

Where $\alpha$ is the step size, $t$ denotes the 
iteration number, $\mathcal{S} = \{\delta : 
\|\delta\|_\infty \leq \epsilon\}$ is the allowed 
perturbation set, and $\Pi_{x+\mathcal{S}}(\cdot)$ 
denotes the projection operator that ensures the 
adversarial example remains within the 
$\epsilon$-ball around the original input $x$. 
PGD is considered a stronger attack than FGSM as 
it iteratively refines the perturbation, making 
it more effective at finding adversarial examples 
that exploit memorized features in the model.